\subsubsection{Bình phương tối thiểu có trọng số (Weighted Least Squares, WLS)}

Phương pháp Bình phương tối thiểu có trọng số (WLS) mở rộng từ OLS nhằm xử lý
trường hợp phương sai của sai số không đồng nhất (heteroskedasticity), tức là:
\begin{equation}
  \mathrm{Var}(\varepsilon_i) = \sigma_i^2.
\end{equation}

Ý tưởng chính là gán trọng số khác nhau cho từng quan sát, trong đó các quan
sát có phương sai lớn sẽ có trọng số nhỏ hơn. Hàm lỗi của WLS được viết lại như
sau: \begin{equation}
  E_D(\boldsymbol{w}) = \frac{1}{2} (\boldsymbol{t} -
  \boldsymbol{\Phi}\boldsymbol{w})^\intercal \boldsymbol{W} (\boldsymbol{t} -
  \boldsymbol{\Phi}\boldsymbol{w}),
\end{equation} trong đó $\boldsymbol{W} = \mathrm{diag}(w_1, \dots, w_N)$ là ma
trận trọng số, với $w_i = 1/\sigma_i^2$.

\paragraph{Phương trình chuẩn cho WLS}

Tương tự OLS, nghiệm tối ưu của bài toán trên được xác định bằng cách giải
phương trình: \begin{equation}
  \boldsymbol{w}^\ast =
  (\boldsymbol{\Phi}^\intercal \boldsymbol{W} \boldsymbol{\Phi})^{-1}
  \boldsymbol{\Phi}^\intercal \boldsymbol{W} \boldsymbol{t}.
\end{equation}

Trong thực tế, để tránh vấn đề ma trận suy biến, ta có thể sử dụng giả nghịch
đảo: \begin{equation}
  \boldsymbol{w}^\ast =
  (\boldsymbol{\Phi}^\intercal \boldsymbol{W} \boldsymbol{\Phi})^\dagger
  \boldsymbol{\Phi}^\intercal \boldsymbol{W} \boldsymbol{t}.
\end{equation}

\paragraph{Ước lượng khả thi bằng FGLS (Feasible Generalized Least
Squares)}

Trong thực tế, các phương sai $\sigma_i^2$ thường không biết trước, do đó ma
trận trọng số $\boldsymbol{W}$ cũng chưa xác định. Phương pháp FGLS được sử
dụng để ước lượng các trọng số này từ dữ liệu.

\subparagraph{Bước 1: Ước lượng ban đầu bằng OLS}

Ước lượng nghiệm ban đầu: \begin{equation}
  \hat{\boldsymbol{w}}_{OLS} =
  (\boldsymbol{\Phi}^\intercal \boldsymbol{\Phi})^{-1}
  \boldsymbol{\Phi}^\intercal \boldsymbol{t}.
\end{equation} Sau đó, tính phần dư: \begin{equation}
  \hat{\varepsilon}_i = t_i - \boldsymbol{\phi}_i^\intercal
  \hat{\boldsymbol{w}}_{OLS}.
\end{equation}

\subparagraph{Bước 2: Mô hình hóa phương sai}

Giả sử phương sai phụ thuộc vào đặc trưng theo công thức: \begin{equation}
  \sigma_i^2 = \exp(\boldsymbol{\phi}_i^\intercal \boldsymbol{\gamma}).
\end{equation} Khi đó, ta thực hiện hồi quy phụ: \begin{equation}
  \log(\hat{\varepsilon}_i^2) = \boldsymbol{\phi}_i^\intercal
  \boldsymbol{\gamma} + u_i,
\end{equation} để ước lượng $\boldsymbol{\gamma}$.

\subparagraph{Bước 3: Ước lượng trọng số}

Từ đó, ta có: \begin{equation}
  \hat{\sigma}_i^2 = \exp(\boldsymbol{\phi}_i^\intercal
  \hat{\boldsymbol{\gamma}}),  \quad \hat{w}_i = \frac{1}{\hat{\sigma}_i^2}.
\end{equation} Xây dựng ma trận trọng số: \begin{equation}
  \hat{\boldsymbol{W}} = \mathrm{diag}(\hat{w}_1, \dots, \hat{w}_N).
\end{equation}

\subparagraph{Bước 4: Thực hiện WLS}

Cuối cùng, nghiệm FGLS được xác định: \begin{equation}
  \hat{\boldsymbol{w}}_{FGLS} =
  (\boldsymbol{\Phi}^\intercal \hat{\boldsymbol{W}} \boldsymbol{\Phi})^{-1}
  \boldsymbol{\Phi}^\intercal \hat{\boldsymbol{W}} \boldsymbol{t}.
\end{equation}

\paragraph{Độ phức tạp tính toán}

Chi phí tính toán của FGLS bao gồm ba bước chính: \begin{itemize}
  \item Ước lượng OLS ban đầu: $\Theta(NM^2 + M^3)$, trong đó $\Theta(NM^2)$ để
    tính $\boldsymbol{\Phi}^\top \boldsymbol{\Phi}$ và $\Theta(M^3)$ để nghịch
    đảo ma trận.

  \item Hồi quy phụ để ước lượng phương sai: $\Theta(NM^2 + M^3)$, do đây cũng
    là một bài toán OLS với cùng số chiều đặc trưng.

  \item Ước lượng WLS: $\Theta(NM^2 + M^3)$, bao gồm việc tính
    $\boldsymbol{\Phi}^\top \boldsymbol{W} \boldsymbol{\Phi}$ và nghịch đảo ma
    trận tương ứng.
\end{itemize} Do đó, tổng chi phí tính toán là $\Theta(NM^2 + M^3)$.  Tương tự
với  $N \gg M$ thì độ phức tạp có thể xấp xỉ $\Theta(NM^2)$, tức là cùng bậc
với OLS.
